% !TEX root = ../memoria.tex
\chapter{Estado del Arte}
\label{cap:capitulo2}

% ============================================================
\section{Evolución de las redes 5G NR}
\label{sec:5g}
% ============================================================

La quinta generación de redes móviles (5G NR, \textit{New Radio}) quedó especificada por
el \textit{3rd Generation Partnership Project} (3GPP) a partir de la Release~15, publicada
en su versión completa en junio de 2018 \cite{3gpp_rel15_2018}. Su diseño rompe con la
filosofía de las generaciones anteriores al separar el plano de control del plano de usuario
(\textit{Control/User Plane Separation}, CUPS) y al definir desde el inicio una arquitectura
de red de acceso radioelétrico (RAN) preparada para dividirse en unidades funcionales
distribuibles: la \textit{Central Unit} (CU), la \textit{Distributed Unit} (DU) y la
\textit{Radio Unit} (RU).

\paragraph{Release 15 (diciembre 2017 -- junio 2018).} Establece la especificación de
referencia para 5G NR. Define dos rangos de frecuencia: FR1 (450~MHz--7,125~GHz,
denominado coloquialmente «sub-6 GHz») y FR2 (24,25--52,6~GHz, o «mmWave»). La interfaz
radio utiliza OFDM con prefijo cíclico tanto en el enlace descendente como en el ascendente,
con subportadoras escalables entre 15~kHz y 240~kHz (\textit{numerology} $\mu = 0$--4).
La Rel.~15 también introduce los primeros despliegues \textit{Non-Standalone} (NSA), en los
que la celda 5G NR comparte el plano de control con una celda LTE anclada, permitiendo
lanzar servicios antes de que exista una core 5G propia.

\paragraph{Release 16 (junio 2020).} Completa la especificación \textit{Standalone} (SA),
donde el núcleo de red es íntegramente 5G~Core (5GC). Introduce soporte completo para
URLLC (\textit{Ultra-Reliable Low-Latency Communications}), V2X (\textit{Vehicle-to-Everything})
y \textit{positioning} sub-métrico mediante señales de referencia de posicionamiento (PRS).
La fiabilidad de enlace (BLER $\leq 10^{-5}$) y la latencia extremo a extremo ($\leq 1$~ms
en el aire) que habilita esta release son las que permiten pensar en gemelos digitales
industriales sincronizados en tiempo real \cite{3gpp_rel16_2020}.

\paragraph{Release 17 (marzo 2022).} Incorpora 5G para IoT masivo a bajo coste
(\textit{Reduced Capability}, RedCap), que reduce la complejidad del terminal a cambio de
limitar el ancho de banda a 20~MHz en FR1. También añade soporte para redes no terrestres
(NTN, \textit{Non-Terrestrial Networks}): satélites en órbita baja y plataformas
estratosféricas (HAPS), que son de especial relevancia para escenarios tácticos en zonas
sin infraestructura fija \cite{3gpp_rel17_2022}.

\paragraph{Escenarios de uso y KPIs.} El estándar IMT-2020 de la ITU-R agrupa los servicios
5G en tres familias: eMBB (\textit{enhanced Mobile Broadband}, pico $> 20$~Gbps en DL),
URLLC (latencia < 1~ms, disponibilidad 99,999~\%) y mMTC (\textit{massive Machine-Type
Communications}, 10$^6$ dispositivos/km$^2$). Las métricas de calidad de señal en la
interfaz radio más habituales son el RSRP (\textit{Reference Signal Received Power})
definido en 3GPP TS 38.215 \cite{3gpp_ts38215_2022}, el RSRQ (\textit{Reference Signal
Received Quality}) y el SINR (\textit{Signal-to-Interference-plus-Noise Ratio}). El RSRP
es la magnitud primaria que produce la simulación de este TFG.

% ============================================================
\section{Redes tácticas y burbujas 5G desplegables}
\label{sec:burbujas}
% ============================================================

Una red táctica 5G es una celda temporal de cobertura reducida que se despliega en minutos
sobre un vehículo, un \textit{flycase} portátil o un dron, pensada para escenarios de
emergencia, operaciones militares o eventos masivos sin cobertura comercial suficiente.
El operador del despliegue controla completamente la core y el acceso radio, lo que elimina
la dependencia de la red comercial y permite priorizar el tráfico de misión crítica
\cite{camp2020tactical5g}.

Las soluciones comerciales más extendidas en el mercado de defensa y protección civil
incluyen la \textit{Mobile Edge Computing Box} de Ericsson (basada en gNodeB de
interior con core 5G integrada en un rack de 19"), la plataforma de Huawei para
redes privadas 5G y el sistema TAC-5G de L3Harris. En España, el programa RAPAZ del
Ejército de Tierra explora el uso de burbujas 4G/5G privadas sobre vehículos tácticos
ligeros. Todos estos sistemas enfrentan el mismo problema: sin una herramienta de
planificación rápida y precisa, el operador táctico no puede estimar la cobertura antes
de posicionar la antena.

% ============================================================
\section{Modelos de propagación radioeléctrica}
\label{sec:propagacion}
% ============================================================

La selección del modelo de propagación es la decisión técnica de mayor impacto en la
precisión de cualquier herramienta de planificación de cobertura. A continuación
revisamos los modelos más relevantes para el rango FR1 (sub-6 GHz).

\subsection{Pérdidas en el espacio libre (FSPL)}
\label{subsec:fspl}

La pérdida en espacio libre es la referencia teórica de mínimas pérdidas, asumiendo que
no existe ningún obstáculo entre TX y RX y que la propagación sigue únicamente la ley del
inverso del cuadrado de la distancia. La fórmula de Friis da:

\begin{equation}
\text{FSPL}~[\text{dB}] = 20\log_{10}(d) + 20\log_{10}(f_c) + 20\log_{10}\!\left(\frac{4\pi}{c}\right)
\label{eq:fspl_cap2}
\end{equation}

donde $d$ es la distancia en metros, $f_c$ la frecuencia en Hz y $c$ la velocidad de la
luz. A 3,5~GHz, por ejemplo, la FSPL a 1~km es aproximadamente 123~dB. Este valor sirve
como cota inferior de las pérdidas reales: cualquier modelo de propagación en exteriores
producirá pérdidas mayores porque incorpora el efecto del terreno, los edificios y la
difracción.

\subsection{Modelos Okumura-Hata y COST-231}
\label{subsec:hata}

El modelo de Okumura \cite{okumura1968field} fue el primer modelo empírico de gran
aceptación para redes celulares. Se construyó a partir de extensas campañas de medida en
Tokio a 200--1920~MHz. Hata \cite{hata1980empirical} lo transformó en ecuaciones analíticas
cerradas, válidas para 150--1500~MHz, que se convirtieron en la norma ITU-R P.529.

El proyecto COST-231 extendió el modelo de Hata hasta los 2000~MHz para dar cobertura a
la banda GSM-1800 y UMTS, añadiendo la constante $C_m$ que diferencia el entorno urbano
grande ($C_m = 3$~dB) del medio \cite{cost231finalreport1999}. La expresión resultante
(ecuación \ref{eq:cost231} del capítulo de implementación) es rápida y no requiere datos
de terreno, lo que la hizo estándar en la planificación de redes 2G/3G. Sin embargo,
su rango de validez estricto termina en 2~GHz, y su extrapolación a 3,5~GHz introduce
errores de 2--5~dB respecto a medidas reales \cite{3gpp_tr38901_2022}.

\subsection{Modelo Longley-Rice (ITM)}
\label{subsec:longleyrice}

El \textit{Irregular Terrain Model} (ITM), desarrollado por Longley y Rice en 1968 en
el ITS de la NTIA (Boulder, Colorado) \cite{longley1968transmission}, es el único de los
modelos considerados que calcula las pérdidas a partir del perfil de elevación real del
terreno entre TX y RX. Para cada enlace, el modelo obtiene el perfil de alturas con el
mismo paso que la cuadrícula de receptores, lo somete a un análisis de difracción múltiple
(método de Fresnel extendido) y estima la contribución de la troposfera a través del
parámetro de irregularidad del terreno $\Delta h$. El rango de validez es 20~MHz a
20~GHz, lo que lo hace directamente aplicable a la banda n78 (3,5~GHz).

A diferencia de los modelos estocásticos de 3GPP, el ITM es determinista: dado el mismo
perfil de terreno, dos simulaciones producen el mismo resultado. Esto es una ventaja para
la planificación táctica, donde el operador necesita saber con certeza si un punto concreto
del mapa tendrá cobertura, no solo la media estadística de la zona.

El ITM fue revisado y validado en el informe NTIA 82-100 \cite{hufford1982guide} y sigue
siendo la referencia del gobierno federal de Estados Unidos para análisis de interferencias
(FCC 47 CFR Part 73). En MATLAB está disponible a partir de R2019b mediante
\texttt{propagationModel('longley-rice')} dentro del \textit{Communications Toolbox}.

\subsection{3GPP TR 38.901}
\label{subsec:tr38901}

El informe técnico 3GPP TR 38.901 \cite{3gpp_tr38901_2022} define los modelos de canal
estocástico para 5G NR en el rango 0,5--100~GHz. A diferencia de los modelos anteriores,
no produce una pérdida de trayecto única sino una realización estadística de canal completa,
con \textit{clusters} de \textit{rays} caracterizados por su retardo, ángulo de salida y
llegada y amplitud. Los escenarios definidos son UMi (\textit{Urban Micro}), UMa
(\textit{Urban Macro}), RMa (\textit{Rural Macro}) y \textit{Indoor Hotspot}.

Para la herramienta de este TFG no utilizamos TR 38.901 directamente como modelo de
propagación porque sus realizaciones estocásticas no son comparables punto a punto con
medidas de campo sin promediar sobre muchos snapshots. Sin embargo, sí empleamos sus
umbrales de RSRP para clasificar la calidad de cobertura (véase el
Cuadro~\ref{tab:rsrp_umbrales} en §4.1.6 del capítulo de implementación).

\subsection{Comparativa de modelos}
\label{subsec:comparativa_modelos}

El Cuadro~\ref{tab:comparativa_modelos} sintetiza las características principales de los
cuatro modelos presentados en esta sección, con el fin de justificar la selección del
modelo Longley-Rice como motor de propagación del gemelo digital.

\begin{table}[htbp]
\centering
\caption{Comparativa de modelos de propagación radioeléctrica.}
\label{tab:comparativa_modelos}
\footnotesize
\begin{tabular}{lcccc}
\hline
\textbf{Criterio} & \textbf{FSPL} & \textbf{COST-231} & \textbf{Longley-Rice} & \textbf{TR 38.901} \\
\hline
Rango de frecuencia      & Cualquiera     & 1,5--2 GHz       & 20 MHz--20 GHz       & 0,5--100 GHz \\
Rango de distancia       & Cualquiera     & 1--20 km         & 1--2000 km           & $<$5 km \\
Datos de terreno         & No             & No               & Sí (SRTM)            & Estadístico \\
Tipo de resultado        & Determinista   & Determinista     & Determinista         & Estocástico \\
Error típico a 3,5 GHz   & 0 dB (teórico) & $\pm$15--20 dB   & $\pm$6--10 dB        & $\pm$4--8 dB \\
Coste computacional      & Mínimo         & Bajo             & Medio                & Alto \\
Uso en este TFG          & Referencia     & Fallback         & \textbf{Principal}   & Validación \\
\hline
\end{tabular}
\end{table}

% ============================================================
\section{Herramientas de planificación radioeléctrica}
\label{sec:herramientas}
% ============================================================

Existen varias herramientas de planificación de cobertura en el mercado. Las evaluamos
desde la perspectiva de su adecuación a un despliegue táctico rápido, donde el operador
necesita resultados en menos de cinco minutos con datos de terreno actualizados.

\paragraph{Xirio Online.} Plataforma SaaS española de INDRA que ofrece cálculo de
cobertura por \textit{ray-tracing} y modelos estadísticos sobre cartografía OSM. Su
precisión es elevada, especialmente en entornos urbanos densos, pero el acceso requiere
suscripción ($\sim$5000~€/año) y el cálculo de un mapa de 100~km$^2$ puede tardar varios
minutos en la capa de servidor compartida. No dispone de API pública documentada para
integración automatizada \cite{xirio2024}.

\paragraph{Atoll (Forsk).} Herramienta de escritorio y de red que soporta LTE, 5G NR y
Wi-Fi. Es el estándar de facto en los departamentos de ingeniería de los operadores europeos.
Gestiona propagación con modelos SPM (\textit{Standard Propagation Model}), Longley-Rice y
\textit{ray-tracing}. El coste de licencia (decenas de miles de euros por sede) y la
dependencia de Windows la hacen inviable para un TFG.

\paragraph{SPLAT!} Herramienta de código abierto para Linux desarrollada por John
Magliacane (KD2BD) \cite{splat2011}, que aplica el modelo Longley-Rice sobre ficheros de
terreno SRTM. Permite calcular curvas de cobertura y enlaces punto a punto con resolución
de 1 arco-segundo (30~m). Su interfaz de línea de comandos es sencilla pero no ofrece
visualización interactiva ni exportación JSON directa, lo que dificulta su integración en
un pipeline moderno.

\paragraph{CloudRF.} Servicio \textit{freemium} web que expone el motor SPLAT!/ITM a
través de una API REST y un \textit{front-end} de mapa. La capa gratuita permite
10~cálculos/día con resolución degradada; la versión de pago soporta alta resolución y
\textit{bulk jobs}. La dependencia de conectividad con los servidores del proveedor la
descarta para uso táctico en zonas sin acceso a internet \cite{cloudrf2024}.

\paragraph{RadioMobile.} Herramienta gratuita de Roger Coude (VE2DBE) para Windows,
orientada a radioaficionados y emergencias. Usa modelos ITM e ITU-R P.370. Su interfaz
ha quedado obsoleta y carece de soporte para 5G NR.

\subparagraph{Gap identificado.} Ninguna de las herramientas anteriores cumple
simultáneamente los tres requisitos clave de un entorno táctico desplegable:
(1)~gratuita y sin dependencia de servidores externos; (2)~exportación estandarizada de
resultados (JSON) para integración con sistemas de mando y control; y (3)~interfaz web
moderna usable por un operador sin formación en RF. Este TFG cubre ese hueco.

% ============================================================
\section{Gemelo digital: concepto y aplicaciones en redes}
\label{sec:gemelo}
% ============================================================

El concepto de \textit{digital twin} (gemelo digital) fue formulado por Michael Grieves
en 2002 en el contexto de la gestión del ciclo de vida del producto (PLM) en la industria
aeroespacial \cite{grieves2014digital}. La idea central es la de un par de entidades: el
sistema físico real y su contraparte virtual, sincronizadas mediante un flujo de datos
bidireccional. La réplica virtual permite simular el comportamiento del sistema en
condiciones que sería costoso o peligroso reproducir en el mundo real, anticipar fallos
y optimizar parámetros en tiempo prácticamente real.

Fuller et al. \cite{fuller2020digital} ofrecen la revisión sistemática más citada del
estado del arte en gemelos digitales, con más de 2000 referencias. Distinguen tres tipos
de gemelo según su nivel de fidelidad:
\begin{itemize}
  \item \textbf{Descriptivo}: réplica estática del sistema en un instante dado
    (modelo CAD, plano de planta).
  \item \textbf{Predictivo}: incluye un modelo dinámico que permite simular el
    comportamiento futuro (simulación de fluidos, electromagnética).
  \item \textbf{Prescriptivo}: cierra el bucle entre simulación y actuación, de forma
    que el gemelo propone (o ejecuta) acciones sobre el sistema real para alcanzar un
    objetivo definido.
\end{itemize}

En el ámbito de las telecomunicaciones, los gemelos digitales han encontrado aplicación en
dos frentes principales. En la \textbf{optimización de redes}: operadores como Ericsson y
Nokia han presentado plataformas de gemelo digital de red que agregan trazas de OSS/BSS y
KPIs de celda para detectar anomalías y sugerir resintonizaciones automáticas de parámetros
RAN \cite{itu2021dt_networks}. En la \textbf{planificación de despliegue}: grupos de
investigación como el de la Universidad de Oulu han propuesto gemelos digitales de zona
urbana para seleccionar emplazamientos de small cells con validación antes del despliegue
físico.

El gemelo digital de este TFG es de tipo \textbf{predictivo}: toma los parámetros de un
gNodeB (frecuencia, potencia, emplazamiento, datos SRTM) y genera un mapa de RSRP
simulado que predice la cobertura real antes de posicionar la antena física. La interfaz
Angular visualiza esa predicción en tiempo real y permite iterar sobre los parámetros de
configuración sin necesidad de un despliegue físico, lo que reduce drásticamente el tiempo
de planificación táctica.

% ============================================================
\section{Trabajos relacionados}
\label{sec:relacionados}
% ============================================================

Varios grupos académicos han abordado problemas parcialmente solapados con el de este TFG.

Nguyen et al. \cite{nguyen2021digital} proponen un gemelo digital de red 5G basado en
simulación de canal con NS-3, con enfoque en la optimización de \textit{beamforming}
masivo (mMIMO). Su plataforma es muy precisa en entornos estáticos conocidos pero no
aborda el problema de la planificación rápida sobre terreno desconocido.

Alraih et al. \cite{alraih2022survey} realizan una revisión de modelos de propagación para
bandas superiores a 6~GHz (\textit{millimeter-wave}) relevantes para 5G y mmWave. Si bien
este TFG se centra en FR1 (sub-6~GHz), algunos de sus resultados sobre la sensibilidad del
error de predicción al modelo elegido son directamente aplicables.

En el ámbito de herramientas abiertas, Gast et al. describen el uso de SPLAT! para
planificación de redes de emergencia 802.11 sobre SRTM y validan el modelo frente a
medidas GPS en campo \cite{gast2005wifi}. Su metodología de validación (RMSE de RSSI
frente a distancia) es la referencia directa para el protocolo de validación experimental
del modelo propuesto en este TFG (Cap.~\ref{cap:capitulo5}).

\paragraph{Diferenciación de este TFG.} Frente a los trabajos anteriores, la contribución
de este proyecto combina cuatro elementos que no aparecen juntos en la literatura revisada:
(1)~integración del modelo Longley-Rice con datos SRTM de 30~m en un script MATLAB
completamente reproducible; (2)~API REST Spring Boot como capa de servicio; (3)~visualización
Angular con heatmap interactivo y selección de punto en el mapa; y (4)~foco explícito en
el escenario de burbuja táctica desplegable, con especial atención a la usabilidad del
operador en campo. El código es abierto y reproducible bajo licencia MIT.
